\section{Limitations, Implications, and Conclusion}
\label{sec:limitations}

\paragraph{What is established.}
Candidate-indexed mass can make a shared-label active experiment depend on roster representation. Exact-quality T2-A and fixed-query equality establish the complete refinement-to-root-cost causal chain, and learned raw-input adapters show that held-out-reference lookup is not necessary for a terminal effect. The missing-frame theorem explains why responses cannot generally recover the intended unit, while Propositions~\ref{prop:factorization} and~\ref{prop:approximate-stability} give exact and approximate design obligations.

\paragraph{Transfer and magnitude boundary.}
The natural anchor is one pair and causes no reliable terminal harm. T2-A uses public or independently acquired references; T1 passes $3/8$, LLM Selector $2/6$, and untargeted additions are mostly benign or beneficial. Most learned same-$s$ tasks are terminal-null or stop. IMDB is disjoint held-out evidence, but its condition was developed on 12,000 rows of the same benchmark. A prospective Yelp audit is statistically positive but below its $.5$-point materiality gate. Zero-utility CIFAR refinements change most paths and satisfy the causal controls, yet also remain submaterial. We therefore establish selective evidence-frame dependence and targeted harm, not a universal relationship between path divergence and practically large terminal loss.

\paragraph{Operational scope.}
The adapters are executable raw-input endpoints, but no endpoint is admitted to a live registry. We do not measure provider incentives, prevalence, secret-task success, or production compromise, and the learned variants share public base architectures rather than independent full-model training. The accountable root is also claim-relative: checkpoint lineage, provider identity, or governance evidence may be needed to authenticate it. A prediction-only distance can merge independent models or split one lineage, as our defense trade-offs illustrate.

\paragraph{Implication.}
The operational requirement is not a universal deduplication metric. An evaluator must declare what a root means for its claim, bind entries to that unit using appropriate provenance or governance evidence, assign fixed root mass, and specify a refinement-invariant within-root summary. When authentication is unavailable, multiplicity sensitivity is part of the uncertainty of the reported selection result and should be audited explicitly.

\paragraph{Conclusion.}
Candidate lists are not innocuous inputs to every active evaluation: under entry-indexed mass, they can encode an unreported prior over accountable alternatives and thereby change the evidence the experiment ever observes. Our theory, exact-quality interventions, and fixed-query controls identify this channel; the retained negative transfers delimit its practical reach. The resulting contribution is an evaluation-unit and invariance requirement, not a claim that candidate multiplication is generally damaging.
